\documentstyle[12pt,newsymb,tabular]{article}
\newtheorem{conjecture}{Conjecture}%[section]
\newtheorem{example}{Example}%[section]
\newtheorem{corollary}{Corollary}%[section]
\newtheorem{lemma}{Lemma}%[section]
\newtheorem{theorem}{Theorem}%[section]
\newtheorem{proposition}{Proposition}%[section]
\newtheorem{definition}{Definition}%[section]
\newtheorem{notations}{Notations}%[section]
\newtheorem{assumption}{Assumption}%[section]
\newtheorem{remark}{Remark}%[section]
\newtheorem{problem}{Problem}%[section]

\begin{document}




\title{Sharp bounds on the $H_p$ means of the derivative of a convex
function for $p=-1$.}
\author{Roger W.\ Barnard\\ and\\  Kent Pearce\\
Texas Tech University\\[2ex]
Dedicated to the memory of our friend Glenn Schober.}

\date{}
\maketitle

\section{Introduction}

For $d>0$ let $D_d=\{z:|z|<d\}$ with $D_1=D$ and let $\partial D_d$
denote the boundary of $D_d$.  Let $S$ be the standard class of
analytic, univalent functions $f$ on $D$, normalized by $f(0)=0$ and
$f'(0)=1$.  Let $K$ denote the well known class of convex functions
in $S$ and let 
\begin{equation}
\label{eq:1}
K_d=\{f\in K: \min_{z\in D} \left|\frac{f(z)}{z}\right|=d\}.
\end{equation}

We will frequently (verbally) identify an analytic, univalent
function $f$ on $D$ with normalization
\renewcommand{\theequation}{\fnsymbol{equation}}
\setcounter{equation}{0}
\begin{equation}
\label{eq:*}
f(0)=0\ {\rm and}\ f'(0)>0
\end{equation}
\renewcommand{\theequation}{\arabic{equation}}
\setcounter{equation}{1}

\noindent with its range $f(D)$ and conversely, since the Riemann mapping
theorem guarantees that we can do so without ambiguity.
Specifically, we will refer to convex domains which contain the
origin and mean the convex functions $f$ which map onto those domains
and satisfy (\ref{eq:*}).

A problem which arose out the authors' work in the early 80's on
omitted value problems for univalent functions, see \cite{2} and
\cite{4}, is the following:  Given $d, \frac{1}{2}, \leq d\leq 1$,
determine a sharp constant $A=A(d)$ such that for any $f\in K_d$
\begin{equation}
\label{eq:2}
I_{-1}(f') =\frac{1}{2\pi} \int\limits^{2\pi}_{0}
\left|\frac{1}{f'(e^{i\theta})}\right| d\theta \leq \frac{A}{d}.
\end{equation}
It follows fairly easily from subordination theory that $A\leq
4/\pi$, but this is not sharp for convex functions.  For $0\leq
\alpha < 1$ let $S^\ast(\alpha)$ denote the usual subclass of $S$ of
starlike functions of order $\alpha$, i.e., a function $f\in S^\ast
(\alpha)$ if and only if $f$ satisfies Re $zf'(z)/f(z)>\alpha$ on
$D$, see \cite{9}.  We will show that $4/\pi$ is sharp for the class
$S^\ast (\frac{1}{2})$ which strictly contains $K$.  In Theorem 3 we
determine for each $\alpha, 0\leq \alpha <1, A=A(\alpha)$ in
(\ref{eq:2}) for the class $S^\ast (\alpha)$ and show that it is
sharp.

Considerable numerical evidence suggested to the authors to make the
following conjecture:

\begin{conjecture}
For each $d, \frac{1}{2} \leq d\leq 1,
A=A(d)=1$ in {\rm (\ref{eq:2})} for the class $K_d$ with equality holding
for all domains which are bounded by regular polygons centered at the
origin.
\end{conjecture}

This conjecture was announced at several talks and conferences over
the last decade including the first author's talk on ``Open problems
in complex analysis'' given at the Symposium on the Occasion of the
Proof of the Bieberbach Conjecture at Purdue University in March
1985.  It also appeared as Conjecture 8 in the first author's ``Open
Problems and Conjectures in Complex Analysis'' in \cite{2}.  It was
thought, by many function theorists, that the conjecture would be
easily settled, given the vast literature on convex functions and the
large research base for determining integral mean estimates, see
\cite{6}.

An initial difficulty was the non-applicability of Baernstein's
general circular symmetrization methods.  The non-convexity of the
circularly symmetrized square shows that convexity, unlike univalence
and starlikeness, is not preserved under circular symmetrization.
Although Steiner symmetrization does preserve convexity, see \cite{10},
it did not appear to be helpful for the problem and, indeed, we will
show that an extremal domain need posses no standard symmetry.

A confusing issue, which also arises, is that the integral means
$I_{-1}(f'_n)$ for the standard approximating functions $f_n$ in $K$,
which map $D$ onto $n$-sided convex polygons and which are defined by
$f'_n (z)=\prod^n_{k=1} (1-ze^{i\theta_k})^{-2\alpha_k}$, $0<\alpha
\leq 1, \sum^n_{k=1}\alpha_k=1$, decrease, as was recently shown in
\cite{1}, when the arbitrarily distributed $\theta_k$ are replaced by
uniformly distributed $t_k=k\pi/n$.  The regular polygons, produced by
the uniformly distributed $t_k$, are the conjectured extremal domains.
 The conjecture suggests that multiplication by the minimum modulus
$d$ must overcome this decrease.

We shall verify the conjecture in Theorem 1.  We also obtain, arising
out of the proof, a rather unexpected sufficient condition for
equality to occur in (\ref{eq:2}) for the classes $K_d$.
Additionally, modifying the proof of Theorem 1, we obtain in Theorem
\ref{th:1*} and its corollary sharp upper and lower bounds for the
integrals means $I_{-1} (f')$.

\section{Statements of the Main Theorems}

Let $K_d$ be defined as in (\ref{eq:1}).  We will say that a convex
curve $\Gamma$ {\em circumscribes} a circle $C$ if the left- and right-hand
tangents at each point of $\Gamma$ are tangent to $C$.

\begin{theorem}\label{th:1}
 Let $d$ be given, $\frac{1}{2}\leq d\leq 1$, and let
$f\in K_d$.  Then
\begin{equation}
\label{eq:3}
\frac{1}{2\pi} \int\limits^{2\pi}_{0}
\left|\frac{1}{f'(e^{i\theta})}\right|d\theta\leq \frac{1}{d}
\end{equation}
with equality holding in {\rm (\ref{eq:3})} if the boundary of $f(D)$
circumscribes $\partial D_d$.
\end{theorem}

Using Theorem 1 and earlier work of the first author we also prove
the following result, which gives a bound for the convex case
of Brennan's conjecture \cite{5} for arbitrary univalent functions.
Let $d$ be given, $\frac{1}{2} \leq d\leq 1$, and let $F_d$ be the
uniquely defined function in $K_d$ which maps $D$ onto the convex
domain having as its boundary an arc, $C_d$, on $\partial D_d$ which
is symmetric about the positive real axis and also
has as its boundary two lines, tangent to $\partial D_d$ at the
endpoints of $C_d$.  Note that $F_1(z)=z$ and
$F_{\frac{1}{2}}(z)=z/(1+z)$.  Note, also, the $\partial F_d(D)$
circumscribes $\partial D_d$.


\begin{theorem}\label{th:2}
Let $d$ be given, $\frac{1}{2}\leq d\leq 1$, and let $f\in K_d$.
Then
\begin{equation}
\label{eq:4}
\displaystyle{\int\hspace{-6pt}\int\hspace{-18pt}\lower
15pt\hbox{$\scriptstyle D$}\hspace{12pt}} |f'(z)|^{-1}dxdy \leq 2\pi\int^1_0
\frac{r^2dr}{F_d(r)}.
\end{equation}
\end{theorem}

Finally, we shall prove

\begin{theorem}\label{th:3}
Let $\alpha$ be given, $0\leq \alpha <1$.  If $f\in S^\ast (\alpha)$
and $\displaystyle{\min_{z\in D} |f(z)/z|=d}$, then
\begin{equation}
\label{eq;5}
\frac{1}{2\pi} \int\limits^{2\pi}_0 \left|\frac{1}{f'
(e^{i\theta})}\right| d\theta \leq \frac{A(\alpha)}{d}
\end{equation}
where
\begin{equation}
\label{eq:6}
A(\alpha)=\frac{4}{\pi} \frac{\arctan
\sqrt{2\alpha-1}}{\sqrt{2\alpha-1}}
\end{equation}
and $A(\alpha)$ is sharp for the class $S^\ast(\alpha)$.
\end{theorem}

\section{Proofs of Main Theorems}

{\bf Proof of Theorem 1.}
\vspace{12pt}

Let $d$ be given, $\frac{1}{2} \leq d \leq 1$, and let $f\in K_d$.
The major idea in proving Theorem 1 is to produce, using the
convexity of the domain $\Omega =f(D)$, two varied domains
$\Omega^\ast$ and $\Omega^{\ast\ast}$ which will satisfy the
inequality
\begin{equation}
\label{eq:7}
\Omega^\ast \subseteq \Omega^{\ast\ast}.
\end{equation}

For any simply connected domain $\Lambda$ which contains the origin
let $\mbox{m.r.}(\Lambda)$ denote the mapping radius of $\Lambda$.  Recall
that if $\Lambda = g(D)$, where $g$ is analytic and univalent on $D$
and $g$ satisfies the normalization (\ref{eq:*}), then
$\mbox{m.r.}(\Lambda)=g'(0)$. Let $\Lambda^+$ be a varied domain of
$\Lambda$ which contains $\Lambda$.  Denote the change in the mapping
radius from the domain $\Lambda$ to $\Lambda^+$ by $\Delta\mbox{m.r.}
(\Lambda^+, \Lambda)$.

The domain containment in (\ref{eq:7}) and subordination will imply
that 
\begin{equation}
\label{eq:8}
\Delta\mbox{m.r.}(\Omega^\ast, \Omega)\leq
\Delta\mbox{m.r.}(\Omega^{\ast\ast}, \Omega)
\end{equation}
from which the conclusion of Theorem \ref{th:1} will be obtained.

The varied domain $\Omega^{\ast\ast}$ is constructed as follows: Let
$\epsilon>0$ be given sufficiently small.  Define
$f^{\ast\ast}_\epsilon$ by $f^{\ast\ast}_\epsilon(z)=(1+\epsilon)
f(z)$ and set $\Omega^{\ast\ast}=f^{\ast\ast}_\epsilon (D)$.  This
has the effect of radially projecting each point $\omega\in \partial
f(D)$ to the point $\omega^{\ast\ast}=(1+\epsilon)\omega$.  This gives
that 
\begin{equation}
\label{eq:9}
[f^{\ast\ast}_\epsilon(z)]'|_{z=0} =\mbox{m.r.} (\Omega^{\ast\ast})
=1+\epsilon\ \ \lower 15pt\hbox{.}
\end{equation}
$$\Delta\mbox{m.r.} (\Omega^{\ast\ast}, \Omega)=\epsilon$$

To construct $\Omega^\ast$ we project each point $\omega\in
\partial\Omega$ outward in the direction of the normal, where it
exists, a distance $\epsilon d$.  Let the extension of $f$ to $\partial D$
also be denoted by $f$.  It is well known that this extension maps
$\partial D$ onto a Jordan curve (on the Riemann sphere if $f$ is
unbounded), $\Gamma=f(\partial D)$, having one-sided tangents
everywhere with the set of points having different one-sided tangents
being at most countable, see \cite{8}.

We define the curve $\Gamma^\ast$, which will be the boundary of
$\Omega^\ast$, as follows: At each point where it exists, let
$n(\omega)$ be the unit outward normal to $\Gamma$ at
$\omega=f(e^{i\theta})$.  Since
$n(\omega)=n[f(e^{i\theta})]=e^{i\theta}f'(e^{i\theta}) /|f'(e^{i\theta})|$,
we will also define at each finite point where the left- and
right-hand tangents differ two limiting normal vectors as
$n^1[f(e^{i\theta})]=n[f(e^{i(\theta-0)})]$ and
$n^2[f(e^{i\theta})]=n[f(e^{i(\theta+0)})]$. We associate with each
point $\omega\in \Gamma$ the point $\omega^\ast =\omega+\epsilon
d[n(\omega)]$ or, where appropriate, the limiting points
$\omega^{\ast j}=\omega + \epsilon d [n^j (\omega)], j=1$ and 2.  To
complete $\Gamma^\ast$ we extend the tangent line to $\Gamma^\ast$ at
each $\omega^{\ast j}$.  Each such tangent line to $\Gamma^\ast$ at
$\omega^{\ast j}$ is parallel to a one-sided tangent line to $\Gamma$
at $\omega$, so that the resulting curve $\Gamma^\ast$ bounds a
convex domain $\Omega^\ast$.  We let $f^\ast_\epsilon$ be the
function which maps $D$ onto $\Omega^\ast $ normalized by
(\ref{eq:*}).

To determine the change in the mapping radius from $\Omega$ to
$\Omega^\ast$ we apply the Hadamard variational formula as developed
in \cite{3} and earlier in \cite{11}.  If $\Gamma$ is bounded, then we
can apply the Hadamard variational formula, cast in the form of the
Julia variation, to obtain
\begin{equation}
\label{eq:10}
\Delta \mbox{m.r.} (\Omega^\ast,\Omega) =\frac{1}{2\pi}
\int\limits^{2\pi}_0 \frac{\epsilon d}{|f'(e^{i\theta})|} d\theta
+o(\epsilon).
\end{equation}

However, if $\Gamma$ is unbounded, we need to show that we can obtain
(\ref{eq:10}) as a limiting argument.  First, we define the varied
domain $\Omega^\ast_0$ by identifying its boundary $\Gamma^\ast_0$.
As in the definition of $\Gamma^\ast$, we associate with each point
$\omega\in\Gamma$ the point $\omega^\ast=\omega+ \epsilon
d[n(\omega)]$ or, where appropriate, the limiting points
$\omega^{\ast j} =\omega + \epsilon d [n^j (\omega)], j=1$ and 2.  To
complete $\Gamma^\ast_0$ we connect $\omega^\ast_1$ to
$\omega^\ast_2$ by the positively oriented circular arc of radius
$\epsilon$ and center $\omega$.  We identify $\Omega^\ast_0$ as the
interior of $\Gamma^\ast_0$ and note that we have $\Omega\subsetneq
\Omega^\ast_0 \subsetneq \Omega^\ast$. 

Let $0<r <1$ and define $f_r$ by $f_r(z)=f((1-r)z)$.  Then, for each
$r$ we have $f_r(D)$ is bounded.  Since $f_r$ is analytic on $\partial
D$, there is for each $\omega_r=f_r(e^{i\theta})\in f_r(\partial D)
=\Gamma_r$ a well-defined outward normal, say, $n(\omega_r)$.  Let
$\omega^\ast_r=\omega_r +\epsilon d [n(\omega_r)]$.  Let
$\Omega^\ast_r$ denote the interior of the curve
$\Gamma^\ast_r=\cup_{\omega_r\epsilon \Gamma_r} \omega^\ast_r$.
Since $\Gamma_r$ is bounded, we can apply the Julia variation to
obtain
$$\Delta \mbox{m.r.} (\Omega^\ast_r, \Omega_r) =\frac{1}{2\pi}
\int\limits^{2\pi}_0 \frac{\epsilon d}{|f'_r (e^{i\theta})|}
d\theta+o(\epsilon).$$
As $r\rightarrow 0, \Omega_r\rightarrow \Omega$ and $\Omega^\ast_r
\rightarrow \Omega^\ast_0$ in the sense of kernel convergence.
Furthermore, $\displaystyle{\frac{1}{|f'_r(e^{i\theta})|} \rightarrow
\frac{1}{|f'(e^{i\theta})|}}$ as $r\rightarrow 0$ except on a
countable set, so we have from the dominated convergence theorem that
$$\lim_{r\rightarrow 0} \int\limits^{2\pi}_0 \frac{1}{|f'_r
(e^{i\theta})|} =\int\limits^{2\pi}_0 \frac{1}{|f'(e^{i\theta})|}
d\theta.$$
It follows then that 
\begin{equation}
\label{eq:11}
\Delta\mbox{m.r.} (\Omega^\ast_0, \Omega) =\frac{1}{2\pi}
\int\limits^{2\pi}_0 \frac{\epsilon d}{|f'(e^{i\theta})|} d\theta
+o(\epsilon).
\end{equation}
It follows from results of the first author in \cite{3} that
$\Delta\mbox{m.r.} (\Omega^\ast, \Omega^\ast_0)=o(\epsilon)$; hence
using (\ref{eq:11}) we have that (\ref{eq:10}) also holds in the case
that $\Gamma$ is unbounded.

To show that (\ref{eq:7}) holds, we will show that for each $\theta
\in [0, 2\pi)$ the radial extension $\omega^{\ast\ast}=(1+\epsilon)
f(e^{i\theta})$ of $\omega =f(e^{i\theta})$ is at least as far from
the origin as the point $\omega_\ast$, the point of intersection of
$\Gamma^\ast =\partial \Omega^\ast$ and the ray $R_t=t\omega: t\geq
1\}$.  See Figure 1.
\vspace*{6in}

\centerline{Figure 1}
\vspace{24pt}

\noindent Let $T_j, j=1$ and 2, denote the left- and
right-hand tangent lines to $\Gamma$ at $\omega$, possibly the same
line, say $T_1$.  Let $z_j$ be the projection of the origin onto
$T_j, l_j$ the line segment joining the origin to $z_j$, and $d_j$
the length of $l_j$.  Let $\theta_j$ be the acute angle between the
ray $R_t$ and the outward normal $n^j(\omega)$ to $\Gamma$ at
$\omega$.  Note that the acute angle between $R_t$ and $l_j$ at the
origin is also equal to $\theta_j$.  Let $t_j$ be the line parallel
to $T_j$ which passes through $\omega^{\ast j} \in \Gamma^\ast$ and
let $x_j$ be the intersection of $t_j$ with $R_t$.  From 
convexity we have
$$|\omega -\omega_\ast| \leq \min_j |\omega-x_j|.$$
We note that when there is a corner at $\omega$, i.e., when $x_1\neq
x_2$, or when $\Gamma$ continues along the tangent line through
$\omega$, then $|\omega-\omega_\ast|=\min_j|\omega-x_j|$, while if
$\Gamma$ does not continue along the tangent line through $\omega$,
then $|\omega-\omega_\ast|<|\omega -x_1|$.

From the similarity of the triangles $[\omega, \omega^{\ast j}, x_j]$
and $[0,z_j,\omega]$ we have
\begin{equation}
\label{eq:12}
\frac{\epsilon d}{|\omega-\omega_\ast|} \geq \max_j \frac{\epsilon
d}{|\omega-x_j|} =\max_j \cos \theta_j =\max_j \frac{d_j}{|\omega|}
\geq \min_j \frac{d_j}{|\omega|} \geq \frac{d}{|\omega|} 
\end{equation}
which gives
\begin{equation}
\label{eq:13}
\frac{\epsilon d}{|\omega-\omega_\ast|} \geq \frac{d}{|\omega|}.
\end{equation}
Thus, $\epsilon|\omega|\geq |\omega-\omega_\ast|$.  Since this is
true for all $\omega=f(e^{i\theta}), \theta \in [0, 2\pi)$, we have
that (\ref{eq:7}) holds.  Combining (\ref{eq:8}), (\ref{eq:9}) and
(\ref{eq:10}) yields
\begin{equation}
\label{eq:14}
\epsilon \geq \frac{\epsilon d}{2\pi} \int\limits^{2\pi}_0
\frac{1}{|f'(e^{i\theta})|} d\theta +o(\epsilon).
\end{equation}
Dividing (\ref{eq:14}) by $\epsilon d$ and letting $\epsilon
\rightarrow 0$ yields the conclusion of the theorem.

We observe that if $\Gamma$ circumscribes $\partial D_d$, then for
each $\omega \in \Gamma$ we have $d=d_1 =d_2, \theta_1=\theta_2$ and
$|\omega-\omega_\ast|=|\omega-x_1|=|\omega-x_2|$, so that equality holds in
(\ref{eq:12}) and (\ref{eq:13}).  In fact, if $\Gamma$ circumscribes
$\partial D_d$, then $\Gamma^\ast=\Gamma^{\ast\ast}$ and $\Omega^\ast
=\Omega^{\ast\ast}$, so that equality holds in (\ref{eq:14}). \hfill
$\blacksquare$
\vspace{12pt}

\noindent {\bf Proof of Theorem 2.}
\vspace{12pt}

Let $F_d$ be the particular convex function, described in the
introduction, which circumscribes $\partial D_d$ and whose boundary
consists of an arc $C_d$ of $\partial D_d$ symmetric about the
positive reals and two lines tangent to $\partial D_d$ at the
endpoints of $C_d$.  Applying Theorem \ref{th:1} in \cite{3}, with
$\alpha =1$, to convex functions $f$ which contain the disk $D_d$
gives the subordination 
\begin{equation}
\label{eq:15}
\log \frac{f(z)}{z} \prec \log \frac{F_d(z)}{z},
\end{equation}
where the subordination in (\ref{eq:15}) means there exists an
analytic function $w$ on $D$ with $|w(z)|\leq |z|$ such that 
$$ \log \frac{f(z)}{z} =\log \frac{F_d (w(z))}{w(z)}.$$
For $d_r=\displaystyle{\min_\theta} |f(r(e^{i\theta}))/r|$ this gives
\begin{equation}
\label{eq:16}
\log d_r \geq \log \left| \frac{F_d(r)}{r}\right|.
\end{equation}
Since, for each $d, \frac{1}{2} \leq d \leq 1, F_d$ is circular
symmetric about the positive reals, it takes is minimum modulus on
$|z|=r$ at $z=r$, $ 0\leq r\leq 1$.  Thus, applying Theorem
\ref{th:1} to $f(rz)/r$, using (\ref{eq:16}) and integrating gives
\begin{equation}
\label{eq:17}
\frac{1}{2\pi} \int\limits^{2\pi}_0 \int\limits^1_0 \frac{r\ dr\
d\theta}{|f' (re^{i\theta})|} \leq \int\limits^1_0 \frac{r\ dr}{d_r}
\leq \int\limits^1_0 \frac{r^2dr}{F_d(r)}.
\end{equation}
From 
(\ref{eq:17}) Theorem \ref{th:2} follows. \hfill $\blacksquare$
\vspace{12pt}

The functions $F_d$ have been shown to be extremal solutions for
growth and subordination problems for the classes $K_d$, which were
studied by Bogucki and Waniurski \cite{bog} and by Barnard and Lewis [2],
respectively.  Also, in a paper related to the authors' earlier work
[1] and [3], Waniurski \cite{wan} has conjectured that the functions $F_d$ are
extremal solutions for an omitted values problem for the classes $K_d$.
Although an explicit closed form for $F_d$ is not known, and appears
to be difficult to obtain,  the recently
developed conformal mapping package CONFPACK \cite{7} can be used to
numerically approximate $F_d$ and, subsequently, the integral in the
right-hand side of (\ref{eq:17}).  We include a table of values.

\newpage

\begin{center}
Table 1\\[2ex]
\begin{tabular}{ccc}
$d$ & $\beta$ & Integral\\
\hline
\\
0.500 & -0.50000 & 0.8333$\bar{3}$ \\
0.525 & -0.47473 & 0.80669 \\
0.550 & -0.44895 & 0.78113\\
0.575 & -0.42268 & 0.75692 \\
0.600 & -0.39589 & 0.73411\\
0.625 & -0.36855 & 0.71263\\
0.650 & -0.34059 & 0.69241\\
0.675 & -0.31192 & 0.67336\\
0.700 & -0.28243 & 0.65540\\
0.725 & -0.25200 & 0.63845\\
0.750 & -0.22046 & 0.62242\\
0.775 & -0.18757 & 0.60724\\
0.800 & -0.15308 & 0.59286\\
0.825 & -0.11658 & 0.57921\\
0.850 & -0.07757 & 0.56796\\
0.875 & & \\
0.900 & & \\
0.925 & 0.06499 & \\
0.950 & 0.12922 & 0.52013\\
0.975 & 0.21574 & 0.50985\\
1.000 & 0.50000 & 0.50000
\end{tabular}
\end{center}


For $\partial F_d(D)$, let $\omega_d$ denote the endpoint of $C_d$ in
the upper half-plane.  The variable $\beta, -\frac{1}{2} \leq \beta
\leq \frac{1}{2}$, in Table 1 parameterizes the family of functions
$F_d, d=d(\beta)$.  If $\beta$ is positive, then the domain
$F_d(D)$ is bounded; otherwise $F_d(D)$ is unbounded.  The argument
of $\omega_d$ is $\pi/2+\beta \pi$.  Also, if $\beta$ is positive,
 then $2\beta\pi$ is
the interior angle for $F_d(D)$ at the finite intersection point of
the lines which bound $F_d(D)$, tangent to $\partial D_d$ at the
endpoints of $C_d$; otherwise $2\beta\pi$ is the interior angle for
$F_d(D)$ at the point of infinity.

The conformal mapping package CONFPACK restricts its applications to
bounded domains.  For $\beta >0$, the domains $F_{d(\beta)}(D)$
were computed directly.  For $\beta<0$ the domains
$F_{d(\beta)}(D)$ were computed by applying a bilinear
transformation to $F_{d(-\beta)}$.  At $\beta=0$ the domain
$F_{d(0)}(D), d(0)\approx 0.8941$, is unbounded and its boundary is
composed of an arc $C_{d(0)}$ of the circle $\partial D_{d(0)}$ and two
lines, parallel to the negative real axis and tangent to the circle
$\partial D_{d(0)}$ at the points of $C_{d(0)}, \ \pm id(0)$. 
The gap in Table 1 occurs because for $\beta$ near $0$ the elongation
of $F_{d(\beta)}(D)$ causes crowding on $\partial D$ of the
pre-images of endpoints of the boundary segments defining
$F_{d(\beta)}(D)$, which causes CONFPACK to fail to converge.  
\vspace{12pt}

\noindent{\bf Proof of Theorem 3}
\vspace{12pt}

It is well known \cite{9} that $f\in S^\ast (\alpha)$ if and only if
\begin{equation}
\label{eq:18}
\frac{zf'(z)}{f(z)} \prec \frac{1+(1-2\alpha)z}{1-z} .
\end{equation}
Thus, applying Littlewood's subordination theorem \cite{6}, we have
for $z=e^{i\theta}$
\begin{eqnarray*}
\frac{1}{2\pi} \int\limits^{2\pi}_0 \frac{1}{|f'(z)|} d\theta
&=&\frac{1}{2\pi} \int\limits_{|z|=1} \left|\frac{f(z)}{zf'(z)} \cdot
\frac{z}{f(z)} \right| |dz| \leq \frac{1}{d} \frac{1}{2\pi}
\int\limits^{2\pi}_0 \left| \frac{f(z)}{zf'(z)}\right| d\theta\\[2ex]
&\leq& \frac{1}{d} \frac{1}{2\pi} \int\limits^{2\pi}_0 \left|
\frac{1-z}{1+(1-2\alpha)z} \right| d\theta =\frac{1}{d} \frac{4}{\pi}
\frac{\arctan \sqrt{2\alpha-1}}{\sqrt{2\alpha-1}}
=\frac{A(\alpha)}{d}.
\end{eqnarray*}
Sharpness follows by using $f_n(z)=z(1-z^n)^{\frac{2(\alpha-1)}{n}}$
and observing that for $d_n=\displaystyle{\min_\theta} |f_n (e^{i\theta})|
=2^{\frac{2(\alpha-1)}{n}}$ we have
$$\lim_{n\rightarrow \infty} \frac{d_n}{2\pi} \int\limits^{2\pi}_0
\left| \frac{1}{f'_n(e^{i\theta})} \right| d\theta =A(\alpha)$$
as claimed.  We note $A(\frac{1}{2}) =4/\pi$.\hfill $\blacksquare$

\section{Note to Theorem 1}

Alternately stated, Theorem \ref{th:1} gives an upper bound for the
integral mean $I_{-1}(f')$ for $f\in K$ in terms of the minimum
distance to the envelope of tangent lines to the boundary of $f(D)$.
An analogous theorem can be proved which gives a lower bound for the
integral mean $I_{-1}(f')$ for $f\in K$ in terms of the maximum
distance to the envelope of tangent lines to the boundary of $f(D)$.
More precisely, in the notation of the proof of Theorem \ref{th:1},
for each $\omega \in \Gamma$, let $T_j, j=1$ and 2, denote the left-
and right-hand tangent lines to $\Gamma$ at $\omega$, possibly the
same line $T_1$.  Let $z_j$ be the projection of the origin on $T_j,
l_j$ the line segment joining the origin to $z_j$, and $d_j$ the
length of $l_j$.  Let $d^\ast
=d^\ast(f)=\displaystyle{\sup_{\omega\in \Gamma} d_j}$.  

We have


\setcounter{section}{1}


\renewcommand{\thetheorem}{{\arabic{section}}{\fnsymbol{theorem}}}
\setcounter{theorem}{0}

\begin{theorem}\label{th:1*}
Let $f\in K$.  Then
\begin{equation}
\label{eq:19}
\frac{1}{d^\ast} \leq \frac{1}{2\pi} \int\limits^{2\pi}_0 \left|
\frac{1}{f'(e^{i\theta})} \right| d\theta
\end{equation}
with equality holding in {\rm (\ref{eq:19})} if the boundary of $f(D)$
circumscribes $\partial D_{d^\ast}$. 
\end{theorem}

\noindent Proof.

The major step in the proof is to show, using the convexity of the
domain $\Omega=f(D)$, there exist two varied domains
$\Omega^{\ast\ast}$ and $\Omega^{\ast\ast\ast}$  which satisfy the
inequality
\begin{equation}
\label{eq:20}
\Omega^{\ast\ast} \subseteq \Omega^{\ast\ast\ast}.
\end{equation}
The domain $\Omega^{\ast\ast}$ is constructed exactly the same as in
Theorem \ref{th:1}, i.e., it is the $(1+\epsilon)$ radial expansion
of $\Omega$.  The domain $\Omega^{\ast\ast\ast}$ is constructed as an
outward normal expansion of $\Omega$, just as $\Omega^\ast$ was, only
we take the constant normal distance to be $\epsilon d^\ast$ instead
of $\epsilon d$.  An analogous argument can be given to show that an
inequality similar to (\ref{eq:12}) holds, only the inequality senses
are all reversed.  In this case, the first inequality in the analog
of (\ref{eq:12}) is more delicate, but it can be verified by
considering a sequence of convex polygonal domains which converge to
$\Omega$ and verifying the inequality for the polygonal domains.
Then, (\ref{eq:20}) holds and a comparison of the mapping radii of
$\Omega^{\ast\ast}$ and $\Omega^{\ast\ast\ast}$ yields 
$$\epsilon \leq \frac{\epsilon d^\ast}{2\pi} \int\limits^{2\pi}_0
\frac{1}{|f' (e^{i\theta})|} d\theta+o (\epsilon)$$
from which Theorem \ref{th:1*} follows. \hfill $\blacksquare$

If we let $d_\ast=d_\ast(f)=\displaystyle{\min_{\omega \in \Gamma}}\ 
d_j$,  then combining both (\ref{eq:3}) and (\ref{eq:19}) we have 
\vspace{12pt}

{\bf Corollary}. {\em  Let $f\in K$.  Then, 
\begin{equation}
\label{eq:21}
\frac{1}{d^\ast} \leq \frac{1}{2\pi} \int\limits^{2\pi}_0 \left|
\frac{1}{f' (e^{i\theta})} \right| d\theta \leq \frac{1}{d_\ast}.
\end{equation}
Equality holds across {\rm (\ref{eq:21})} if for some $d, \frac{1}{2} \leq
d\leq 1, \partial f(D)$ circumscribes $\partial D_d$.}
\vspace{12pt}

{\bf Remark}.  For $f\in K$ let $L=$ length $\partial f(D)$ and let $F$ denote
the inverse of $f$.  The isoperimetric inequaltiy states
$$\frac{2\pi d_\ast}{L} \leq 1.$$
Schwarz's inequality, applied to the integral means in (21) --
rewritten in terms of $F$, yields
\begin{equation}
\label{eq:22}
\frac{2\pi}{L}\leq \frac{1}{2\pi} \int^{2\pi}_0
\left|\frac{1}{f'(e^{i\theta})} \right|d\theta.
\end{equation}
Thus, (21) combined with (\ref{eq:22}) gives an intermediate term to
the isoperimetric inequality.


\begin{thebibliography}{99}



\bibitem{2} Barnard, R.W.\ ``Open problems and conjectures in complex
analysis'', in {\em Computational Methods and Function Theory},  Lecture
Notes in Math.\ No.\ 1435 (Springer-Verlag, 1990), 1-26.

\bibitem{3} Barnard, R.W.\ and Lewis, J.L.\ ``Subordination theorems
for some classes of starlike functions,'' {\em Pacific J.\
Math.}, 56 (1975), 333-366.

\bibitem{4} Barnard, R.W., and Pearce, K.\ ``Rounding corners of
gearlike domains and the omitted area problem'', {\em J.\ Comput.\ Appl.\
Math}. 14 (1986), 217-226.

\bibitem{bog} Bogucki, Z., and Waniurski, J. ``On univalent functions
whose values cover a fixed disk,'' {\em Ann.\ Univ.\ Mariae
Curie-Sklodowska}, Sect.\ A 22 (1968), 39-43.

\bibitem{5} Brennan, J.E. ``The integrability of the derivative in
conformal mapping,'' {\em J.\ London Math.\ Soc.} (2) 18 (1978), 
261-272.

\bibitem{6} Duren, P.\ {\em Univalent Functions},  Springer-Verlag,
New York, 1980.

\bibitem{7} Hough, D.M. ``User's guide to CONFPACK,'' IPS, Research
Report, No.\ 90-11; ETH-Zentrum, Zurich, Switzerland, 1990.

\bibitem{8} Krzyz, J.\ Distortion theorems for bounded convex
functions, II.  {\em Ann.\ Univ.\ Mariae Curie-Sklodowska}, Sect.\ A 14
(1960), 7-18.

\bibitem{9} Robertson, M.S.\ ``On the theory of univalent functions'',
{\em Ann.\ of Math.} 37 (1936), 374-408.

\bibitem{10} Valentine, F.A. {\em Convex Sets}, McGraw-Hill, New York, 1964.

\bibitem{wan} Waniurski, J. ``On values omitted by convex univalent
mappings'', {\em Complex Variables} 8 (1973), 173-180.

\bibitem{11}  Warschawski, S.E. ``On Hadamard's variation formula for
Green's function'' {\em J.\ Math.\ Mech.}\ 9 (1960), 497-512.

\bibitem{1} Zheng, J.\ ``Some extremal problems
involving $n$ points on the unit circle'', Dissertation, 
Washington University (St.\ Louis), 1991.

\end{thebibliography}
\end{document}